%mac
%\loai{Xác định các đại lượng từ phương trình}
\begin{vidu}%[1P1-3.1K][Loai1]
Một xe máy chuyển động dọc theo trục $Ox$ có phương trình chuyển động: $x = 60 - 60t$ (với $x$ đo bằng $km$, $t$ đo bằng $h$).
\begin{vdc}
Xe máy chuyển động theo chiều dương hay chiều âm của trục $Ox$.
\end{vdc}
\begin{vdc}		
Xác định thời điểm xe máy đi qua gốc tọa độ.
\end{vdc}
\loigiai{
\begin{itemize}
\item[$\bullet$] 
Từ phương trình tổng quát: $x=x_0+vt$ so sánh với phương trình đề ra ta có: $x_0 = 60$ km; $v = -60$ km/h.
\item[a)] Vì $v = -60\ km/h < 0$ nên xe máy đang chuyển động theo chiều âm của trục tọa độ.
\item[b)] Khi xe máy đi qua gốc tọa độ thì:
\begin{align*}
x=0\Leftrightarrow 60-60t=0\Rightarrow t=1\ h.
\end{align*}
\end{itemize}
}
\end{vidu}
%\loai{Viết phương trình chuyển động của vật}
\begin{vidu}%[1P1-3.1K][Loai1]
Một học sinh đi xe đạp chuyển động thẳng đều với tốc độ $12\ km/h$ từ $A$ đến $B$ rồi đến $C$. Biết $B$ cách $A$ đoạn $4\ km$ và $C$ cách $B$ đoạn $2\ km$. Viết phương trình chuyển động (phương trình tọa độ) của xe đạp trong các trường hợp sau:
\begin{vdc}
Chọn gốc tọa độ tại $A$, gốc thời gian lúc học sinh xuất phát từ $A$ và chiều dương là chiều chuyển động.
\end{vdc} 
\begin{vdc} 
Chọn gốc tọa độ tại $B$, gốc thời gian lúc học sinh xuất phát từ $A$ và chiều dương là chiều từ $C$ đến $A$.
\end{vdc}
\loigiai{
\begin{itemize}
\item[a)] Toạ độ ban đầu và vận tốc của vật: $x_0 = 0$, $v = 12\ km/h$.\\
Phương trình chuyển động của xe đạp là: $x = 12t$ (x đo bằng km; t đo bằng h).
\item[b)] Toạ độ ban đầu và vận tốc của vật: $x_0 = 4$ km, $ v = -12$ km/h.\\
Phương trình chuyển động của xe đạp là: $x=4-12t$ (x đo bằng km; t đo bằng h).
\end{itemize}	
}
\end{vidu}
%\loai{Hai xe chuyển động ngược chiều, khởi hành cùng lúc}

%\loai{Hai xe chuyển động cùng chiều, khác thời điểm}
\begin{vidu}%[1P1-3.1K][Loai1]
Lúc $7h$ một người đi bộ từ $A$ đến $B$ với tốc độ $v_1 = 4\ km/h$. Lúc $9h$ một người đi xe đạp từ $A$ đuổi theo với tốc độ $v_2 = 12\ km/h$. Chọn trục tọa độ $Ox$ có gốc $O$ trùng với vị trí $A$, chiều dương là chiều từ $A$ đến $B$. Gốc thời gian là lúc $7h$.
\begin{vdc}
Viết phương trình chuyển động của mỗi người.
\end{vdc}
\begin{vdc}
Lúc mấy giờ họ cách nhau $2\ km$.
\end{vdc}
\loigiai{
a) Vì chọn gốc thời gian là lúc 7h nên ta có: $\begin{cases}
\text{Thời gian chuyển động của người đi bộ:} t\\
\text{Thời gian chuyển động của người đi xe đạp:} t-2\\
\end{cases}$
\begin{itemize}
\item Vì gốc tọa độ O trùng với A nên ta có: $\begin{cases}
x_{01}=0.\\
x_{02}=0.\\
\end{cases}$
\item Hai người đi theo chiều dương trục Ox nên $v > 0 \Rightarrow  \begin{cases}
v_1=4\ km/h.\\
v_2=12\ km/h.\\
\end{cases}$
\item Phương trình chuyển động của mỗi người là: $\begin{cases}
x_1=4t.\\
x_2=12\left(t-2\right).\\
\end{cases}$(với x đo bằng km; t đo bằng h).
\end{itemize}
b) Khi họ cách nhau 2 km thì $\left|x_1-x_2\right|=2\Rightarrow x_1-x_2=\pm 2$.
\begin{itemize}
\item \textit{Trường hợp 1:} Khi người đi xe đạp chưa vượt qua người đi bộ
\begin{align*}
x_1-x_2=2\Leftrightarrow 4t-12\left(t-2\right)=2\Rightarrow t=2,75h = 2\ \text{giờ}\ 45\ \text{phút}.
\end{align*}
Vậy người đi xe đạp cách người đi bộ 2 km lần đầu vào lúc 9 giờ 45 phút.
\item \textit{Trường hợp 2:} Khi người đi xe đạp vượt qua người đi bộ
\begin{align*}
x_1-x_2=-2\Leftrightarrow 4t-12\left(t-2\right)=-2\Rightarrow t=3,25h = 3\  \text{giờ}\ 15\ \text{phút}.
\end{align*}
Vậy người đi xe đạp cách người đi bộ 2 km lần thứ 2 vào lúc 10 giờ 15 phút.
\end{itemize}}
\end{vidu}
%\loai{Hai xe chuyển động vuông góc}
\begin{vidu}%[1P1-3.1G][Loai1]
\immini[thm]{Trong hệ tọa độ $xOy$ (hình bên), có hai vật nhỏ $A$ và $B$ chuyển động thẳng đều. Lúc bắt đầu chuyển động, vật $A$ ở $O$ và cách vật $B$ một đoạn $100\ m$. Biết tốc độ của vật $A$ là $v_A = 6\ m/s$ theo hướng $Ox$, tốc độ của vật $B$ là $v_B = 2\ m/s$ theo hướng $Oy$.}{\begin{tikzpicture}[scale=1,thick,>=stealth',draw=Mapcolor]
\tkzDefPoints{0/0/a, 0/-2/b, 0/2/y, 3/0/x}
\draw[->](b)--(y)node[right]{$y$};
\draw[->](a)--(x)node[above]{$x$};
\draw[Mapcolor,very thick,->](b)--++(90:0.7)node[black,right]{$\vec{v}_{B}$};
\draw[Mapcolor,very thick,->](a)--++(0:1.5)node[black,above]{$\vec{v}_{A}$};
\tkzDrawPoints[size=3](a,b)
\node at (a)[left]{$A$};
\node at (b)[left]{$B$};
\node at (a)[below right]{$O$};
\end{tikzpicture}}
\begin{vdc}
 Viết phương trình chuyển động của mỗi vật theo trục tọa độ riêng của chúng.
\end{vdc}
\begin{vdc}
Sau thời gian bao lâu kể từ khi bắt đầu chuyển động, hai vật $A$ và $B$ lại cách nhau $100\ m$.
\end{vdc}
\begin{vdc}
Xác định khoảng cách nhỏ nhất giữa hai vật $A$ và $B$.
\end{vdc}
\haidong{20}
\loigiai{
\immini[thm]{
\begin{itemize}
\item[a)] 	Phương trình chuyển động của vật A: $x_A=6t$.
\begin{itemize}
\item Phương trình chuyển động của vật B: $x_B=-100+2t$.
\end{itemize}
\item[b)] 	Khoảng cách giữa A và B sau t giây: \\$d^2=\left(x_A\right)^2+\left(x_B\right)^2\Leftrightarrow d^2=\left(-100+2t\right)^2+36t^2\left(*\right)$.
\begin{itemize}
\item Khi khoảng cách $d = 100$ m. Ta có: $\begin{cases}
100^2=\left(-100+2t\right)^2+36t^2.\\
\Rightarrow 40t^2-400t=0\Rightarrow t=10\ s.\\
\end{cases}$
\end{itemize}
\end{itemize}
}{\begin{tikzpicture}[scale=1,thick,>=stealth']
\tkzDefPoints{0/0/a, 2.5/0/a1, 0/-3/b, 0/-2/b1, 0/1/y, 3.2/0/x}
\draw[->](b)--(y)node[right]{$y$};
\draw[->](a)--(x)node[above]{$x$};
\draw[blue,very thick,->](b)--++(90:0.7)node[black,right]{$\vec{v}_{B}$};
\draw[red,very thick,->](a)--++(0:1.5)node[black,above]{$\vec{v}_{A}$};
\tkzDrawSegments[dashed](b1,a1)
\tkzDrawPoints[size=3](a,b,a1,b1)
\node at (a)[left]{$A$};
\node at (a1)[below right]{$A_{1}$};
\node at (b1)[left]{$B_{1}$};
\node at (b)[left]{$B$};
\node at (a)[below right]{$O$};
\end{tikzpicture}}
\begin{itemize}
\item[c)] 	Biến đổi (*): $40t^2-400t+100^2=d^2\left(**\right)$. Ta có:
\begin{align*}
\begin{cases}
40t^2-2.\left(2\sqrt{10}t\right).\left(10\sqrt{10}\right)+10^2.10+9000=d^2\left(**\right)\\
\Leftrightarrow {\left(2\sqrt{10}t-10\sqrt{10}\right)}^2+9000=d^2\Rightarrow d_{\min}=30\sqrt{10}\ m.\\
\end{cases}
\end{align*}
\end{itemize}
\begin{note}
\begin{itemize}
\item[$\bullet$] Có thể tính theo biệt thức đenta ở lớp 9 từ phương trình bậc 2 như sau:
\begin{align*}
{\Delta}'={\left(b'\right)}^2-ac=200^2-40\left(100^2-d^2\right)=40d^2-360000.
\end{align*}
\item[$\bullet$] Điều kiện có nghiệm là: ${\Delta}'\ge 0\Leftrightarrow 40d^2-360000\ge 0\Rightarrow d\ge 30\sqrt{10}\Rightarrow d_{\min}=30\sqrt{10}\ m$.
\end{itemize}	
\end{note}	
}
\end{vidu}









%mac
\begin{vidu}
Số liệu về tọa độ và thời gian của chuyển động thẳng của một xe ô tô đồ chơi chạy bằng pin được ghi trong bảng:
\begin{center}
\begin{tabular}{|c|c|c|c|c|c|c|}
\hline 
Tọa độ (m)&1  &3  &5  &7  &7  &7  \\ 
\hline 
Thời gian (s)&0  &1  &2  &3  &4  &5  \\ 
\hline 
\end{tabular}
\end{center}
Dựa vào bảng này để:
\begin{vdc}
Vẽ đồ thị tọa độ - thời gian của chuyển động.
\end{vdc}
\begin{vdc}
Tính vận tốc của xe trong $3\ s$ đầu.
\end{vdc}	 	
\haidong{18}
\loigiai{
\immini[thm]{\begin{itemize}
\item[a)] Vẽ đồ thị tọa độ – thời gian:
\item[b)] Mô tả chuyển động của xe:
\item[-] Từ 0 – 3 giây: xe chuyển động thẳng.
\item[-] Từ giây thứ 3 đến giây thứ 5: xe đứng yên (dừng lại)
\item[c)] Độ dịch chuyển của xe trong 3 giây đầu là: $d=7-1=6\ m$.
\item[-] Vận tốc của xe trong 3 giây đầu là: $v=\dfrac{d}{\Delta t}=\dfrac{6}{3}=2\ m/s$.
\end{itemize}}{\begin{tikzpicture}[scale=1,thick,>=stealth',y=0.5cm]
\tkzDefPoints{0/0/o, 0/-1/i, 5.5/0/t, 0/8/d}
\tkzDrawSegments[very thick,blue,->](o,t i,d)
\tkzLabelPoint[left](o){$O$}
\node[right] at (d){x (m)};
\node[above] at (t){t (s)};
\foreach  \x  [evaluate=\x as \z using int(1*\x)] in  {1,2,...,5}
\node[below] at(\x,0){$\z$};
\foreach  \x  [evaluate=\x as \z using int(1*\x)] in  {1,2,...,5}
\draw(\x,-0.05)--(\x,0.05);
\foreach  \y  [evaluate=\y as \w using int(1*\y)] in  {1,3,...,7}
\node[left] at(0,\y){$\w$};
\foreach  \y  [evaluate=\y as \z using int(1*\y)] in  {1,3,...,7}
\draw(-0.05,\y)--(0.05,\y);
\draw[blue,line width=2pt](0,1)--(3,7)--(5,7);
\draw[dashed](0,7)--(3,7)--(3,0);
\end{tikzpicture}}
}	
\end{vidu}
\begin{vidu}
Hình bên là đồ thị độ dịch chuyển - thời gian của một người đang bơi trong một bể bơi dài $50\ m$. Đồ thị này cho biết những gì về chuyển động của người đó?
\immini{\begin{listEX}
\item Trong $25$ giây đầu mỗi giây người đó bơi được .......... mét.\\ 
Vận tốc của người đó là .......... $m/s$. 
\item Từ giây thứ .......... đến giây thứ .......... người đó không bơi.
\item Từ giây $35$ đến giây $60$ người đó bơi theo chiều .......... 
\item Khi người đó bơi từ $B$ đến $C$ thì độ dịch chuyển có độ lớn là .......... mét và vận tốc của người đó có độ lớn là .......... $m/s$.
\item  Trong cả quá trình bơi thì độ dịch chuyển của người đó là ..........  mét và vận tốc của người đó có độ lớn là .......... $m/s$.
\end{listEX}	}{\begin{tikzpicture}[scale=1,thick,>=stealth',draw=Mapcolor]
\tkzDefPoints{0/0/o, 7/0/t, 0/7/d, 2.5/5/a, 3.5/5/b, 6/2.5/c, 1/2/g, 2.5/2/h}
\draw  [help  lines, xstep=0.5cm,ystep=0.5cm]  (0,0)  grid  (6,6);
\tkzDrawSegments[very thick,->](o,t o,d)
\tkzDrawSegments[very thick,Mapcolor](o,a a,b b,c)
\tkzDrawSegments[very thick,Mapcolor,<->](g,h a,h)
\tkzMarkAngles[size=0.4cm](h,g,a)
\tkzLabelAngles[pos=0.7](h,g,a){$\alpha$}
\tkzLabelSegment[below](g,h){$\Delta t$}
\tkzLabelSegment[right](a,h){$\Delta d$}
\tkzLabelPoint[above left](a){$A$}
\tkzLabelPoint[above right](b){$B$}
\tkzLabelPoint[right](c){$C$}
\tkzLabelPoint[below](t){$t\ (s)$}
\tkzLabelPoint[left](d){$d\ (m)$}
\tkzLabelPoint[below left](o){$O$}
\tkzLabelSegment[sloped,yshift=15pt](o,d){Độ dịch chuyển}
\foreach  \x  [evaluate=\x as \z using int(10*\x)] in  {0.5,1,...,6}
\node[below] at(\x,0){$\z$};
\foreach  \x  [evaluate=\x as \z using int(10*\x)] in  {1,2,...,6}
\node[left] at(0,\x){$\z$};
\end{tikzpicture}}
\haidong{31}
\loigiai{
\begin{itemize}
\item[1.] Từ đồ thị ta thấy, trong 25 s đầu người đó chuyển động thẳng từ O – A và không đổi chiều, độ dịch chuyển trong 25 s đầu là 50 m.
\item[-] Suy ra: Mỗi giây người đó bơi được: $\dfrac{50}{25}=2\ m$.
\item[-] Vận tốc của người đó là: $v=\dfrac{d}{t}=\dfrac{50}{25}=2\ m/s$.
\item[2.] Từ A – B: người đó không bơi $\Rightarrow$ Người đó không bơi từ giây 25 đến giây 35.
\item[3.] Từ giây 35 đến giây 60 người đó bơi ngược chiều dương.
\item[4.] Dựa vào đồ thi:
\item[-] Tại B: $d_1=50\ m;\ t_1=35\ s$.
\item[-] Tại C: $d_2=25\ m;\ t_2=60\ s$.
\item[-] Từ B $\to$ C, độ dịch chuyển là: $\Delta d=d_2-d_1=25-50=-25\ m$.
\item[-] Vận tốc của người đó khi bơi từ B $\to$ C là: $v=\dfrac{\Delta d}{\Delta t}=\dfrac{-25}{60-35}=-1\ m/s$
\item[5.] 
\item[-] Độ dịch chuyển của người đó trong cả quá trình bơi là: $\Delta d=25\ m$
\item[-] Vận tốc của người đó trong cả quá trình bơi là: $v=\dfrac{\Delta d}{\Delta t}=\dfrac{25}{60}=\dfrac{5}{12}\approx 0,417\ m/s$.
\end{itemize}
}
\end{vidu}

\begin{vidu}
\immini[thm]{Đồ thị độ dịch chuyển - thời gian trong chuyển động thẳng của một xe ô tô đồ chơi điều khiển từ xa được vẽ ở hình bên.
}{\begin{tikzpicture}[scale=1,thick,>=stealth',scale=1.3,draw=Mapcolor]
\tkzDefPoints{0/0/o, 5.5/0/t, 0/3/d, 1/2/a, 2/2/b, 4.5/-0.5/c, 5/-0.5/e}
\draw  [help  lines, xstep=0.5cm,ystep=0.5cm]  (0,-1)  grid  (5,2.5);
\tkzDrawSegments[very thick,->](o,t o,d)
\tkzDrawSegments[very thick,Mapcolor](o,a a,b b,c c,e)
\tkzLabelPoint[above ](t){$t\ (s)$}
\tkzLabelPoint[left](d){$d\ (m)$}
\tkzLabelPoint[below left](o){$O$}
\foreach  \x  [evaluate=\x as \z using int(2*\x)] in  {0.5,1,...,5}
\node[below] at(\x,0){$\z$};
\foreach  \x  [evaluate=\x as \z using int(2*\x)] in  {-0.5,0.5,1,...,2.5}
\node[left] at(0,\x){$\z$};
\end{tikzpicture}}
\begin{itemize}
\item[a)] Từ $t=0\ s$ đến $t=2\ s$ xe \dotfill\\
Từ $t=2\ s$ đến $t=4\ s$ xe \dotfill\\
Từ $t=4\ s$ đến $t=8\ s$ xe \dotfill\\
Từ $t=8\ s$ đến $t=9\ s$ xe \dotfill\\
Từ $t=9\ s$ đến $t=10\ s$ xe \dotfill
\item[b)] Vị trí của xe so với điểm xuất phát của xe ở giây thứ $2$ là $d_2=..........$, giây thứ $4$ là  $d_4=..........$, giây thứ $8$ là  $d_8=..........$ và giây thứ $10$ là  $d_{10}=..........$.
\item[c)] Tốc độ của xe trong $2$ giây đầu là $........... m/s$, từ giây $2$ đến giây $4$  là $........... m/s$ và từ giây $4$ đến giây $8$  là $........... m/s$.
\item[d)] Vận tốc của xe trong $2$ giây đầu là $........... m/s$, từ giây $2$ đến giây $4$  là $........... m/s$ và từ giây $4$ đến giây $8$  là $........... m/s$.

\end{itemize}	 
\haidong{20}	
\loigiai{
\begin{itemize}
\item[a)] Mô tả chuyển động của xe:
\item[-] Trong 2 giây đầu: xe chuyển động thẳng
\item[-] Từ giây thứ 2 đến giây thứ 4: xe đứng yên
\item[-] Từ giây thứ 4 đến giây thứ 10: xe chuyển động thẳng theo chiều ngược lại.
\item[-] Từ giây thứ 9 đến giây thứ 10: xe dừng lại.
\item[b)] Ở giây thứ 2: xe ở vị trí cách điểm xuất phát 4 m.
\item[-] Ở giây thứ 4: xe ở vị trí cách điểm xuất phát 4 m
\item[-] Ở giây thứ 8: xe trở về vị trí xuất phát
\item[-] Ở giây thứ 10: xe ở vị trí cách điểm xuất phát 1 m theo chiều âm
\item[c)] Xác định tốc độ và vận tốc của xe:
\item[-] Trong 2 giây đầu, xe chuyển động thẳng, không đổi chiều nên tốc độ bằng vận tốc: $v=\dfrac{d}{t}=2\ m/s$.
\item[-] Từ giây 2 đến giây 4: xe đứng yên nên vận tốc và tốc độ của xe đều bằng 0.
\item[-] Từ giây 4 đến giây 8:
\item[+] Tốc độ: $v=\dfrac{s}{t}=1\ m/s$.
\item[+] Vận tốc: $v=\dfrac{\Delta d}{\Delta t}=\dfrac{0-4}{8-4}=-1\ m/s$.

\end{itemize}
}
\end{vidu}

%\begin{vidu}
%	Hãy vẽ đồ thị độ dịch chuyển - thời gian trong chuyển động của chuyển độngnêu ở bẳng số liệu sau đây:\begin{center}
%		\begin{tabular}{|c|c|c|c|c|c|c|c|}
%			\hline 
%			Độ dịch chuyển (m)&  0&  200&  400&  600&  800&  1000&800  \\ 
%			\hline 
%			Thời gian (s)&  0&  50&  100&  150&  200&  250&300  \\ 
%			\hline 
%		\end{tabular} 
%	\end{center}
%	Vẽ đồ thị: trên trục tung (trục độ dịch chuyển) 1 cm ứng với 200 m; trên trục hoành (trục thời gian) 1 cm ứng với 50 s.
%	\loigiai{
%		\begin{center}
%			\begin{tikzpicture}[scale=1,thick,>=stealth']
%				\tkzDefPoints{0/0/o, 0/-1/i, 6.5/0/t, 0/5.5/d}
%				\tkzDrawSegments[very thick,blue,->](o,t i,d)
%				\tkzLabelPoint[left](o){$O$}
%				\node[right] at (d){d (m)};
%				\node[above] at (t){t (s)};
%				\foreach  \x  [evaluate=\x as \z using int(50*\x)] in  {1,2,...,6}
%				\node[below] at(\x,0){$\z$};
%				\foreach  \x  [evaluate=\x as \z using int(50*\x)] in  {1,2,...,6}
%				\draw(\x,-0.05)--(\x,0.05);
%				\foreach  \y  [evaluate=\y as \w using int(200*\y)] in  {1,2,...,5}
%				\node[left] at(0,\y){$\w$};
%				\foreach  \y  [evaluate=\y as \z using int(200*\y)] in  {1,2,...,5}
%				\draw(-0.05,\y)--(0.05,\y);
%				\draw[dashed](0,1)--(1,1)--(1,0)(0,5)--(5,5)--(5,0)(0,4)--(6,4)--(6,0);
%				\draw[blue,line width=2pt](0,0)--(5,5)--(6,4);
%			\end{tikzpicture}
%		\end{center}
%	}
%\end{vidu}
\begin{vidu}%[1P1K3-1][Loai1]
Lúc $6$ giờ sáng, một người đi xe máy khởi hành từ $A$ chuyển động với tốc độ $36\ km/h$ đi về $B$. Cùng lúc một người đi xe đạp chuyển động với tốc độ  $v_2$ xuất phát từ $B$ đi đến $A$. Khoảng cách $AB = 108\ km$. Hai người gặp nhau lúc $8$ giờ. Tốc độ của xe đạp là
\choice
{$12\ km/h$}
{\True $18\ km/h$}
{$10\ km/h$}
{$22\ km/h$}
\haidong{13}
\loigiai{
\begin{itemize}
\item Hai người gặp nhau lúc 8 giờ, tức là sau 2 giờ chuyển động. 
\item Chọn gốc tọa độ tại A, gốc thời gian là lúc hai xe bắt đầu xuất phát, chiều dương từ A đến B.
\begin{itemize}
\item Phương trình chuyển động của xe 1:
$x_{1}=x_{01}+v_1t=36t$.
\item Phương trình chuyển động của xe 2:
$x_2=x_{02}+v_2t=108-v_2t$.
\end{itemize}
\item Lúc hai xe gặp nhau
\begin{align*}
x_1=x_2\Leftrightarrow 36t=108-v_2t\Leftrightarrow (36+v_2)t=108\tag{1}
\end{align*}
\item Thay \bth{t=2} vào (1) suy ra: \bth{v_2=18\ km/h}. 
\item Vậy tốc độ của xe đạp là 18 km/h.
\end{itemize}
}
\end{vidu}






